\documentclass[conference]{IEEEtran}
\IEEEoverridecommandlockouts
% The preceding line is only needed to identify funding in the first footnote. If that is unneeded, please comment it out.
\usepackage{cite}
\usepackage{amsmath,amssymb,amsfonts}
\usepackage{algorithmic}
\usepackage{graphicx}
\usepackage{textcomp}
\usepackage{xcolor}
\def\BibTeX{{\rm B\kern-.05em{\sc i\kern-.025em b}\kern-.08em
    T\kern-.1667em\lower.7ex\hbox{E}\kern-.125emX}}

\begin{document}

\title{High-Performance Deep Learning Pipeline for Robust OCR Preprocessing and Text Region Extraction}

\author{\IEEEauthorblockN{1\textsuperscript{st} Author Name}
\IEEEauthorblockA{\textit{dept. name of organization (of Aff.)} \\
\textit{name of organization (of Aff.)}\\
City, Country \\
email address or ORCID}
\and
\IEEEauthorblockN{2\textsuperscript{nd} Author Name}
\IEEEauthorblockA{\textit{dept. name of organization (of Aff.)} \\
\textit{name of organization (of Aff.)}\\
City, Country \\
email address or ORCID}
}

\maketitle

\begin{abstract}
Optical Character Recognition (OCR) systems require robust preprocessing to accurately extract text from complex backgrounds, diverse orientations, and distorted perspectives. This paper presents a high-performance, C++-based preprocessing pipeline optimized for text Region of Interest (ROI) extraction. Our system incorporates the ONNX Runtime for deploying DBNet (Differentiable Binarization Network) for initial text detection. To address fragmented detection boxes and varying text shapes, we employ robust geometry processing algorithms, including Stable Component Analysis and polygon manipulation via the Clipper library. Extracted regions undergo perspective unwarping and localized binarization to maximize OCR accuracy. We demonstrate the structural advantages of a compiled C++17 execution environment using the ONNX Intermediate Representation (IR) v9 interface, offering deployment consistency and cross-platform compatibility. 
\end{abstract}

\begin{IEEEkeywords}
OCR, Preprocessing, DBNet, ONNX, Text Detection, Perspective Unwarping, C++
\end{IEEEkeywords}

\section{Introduction}
Optical Character Recognition (OCR) has become an essential component traversing countless enterprise, mobile, and web applications. However, unconstrained physical environments—such as documents scanned with mobile phones, receipts, or natural scene texts—present significant challenges such as irregular lighting conditions, distorted perspectives, and sophisticated background shapes. Before sophisticated recognition engines (e.g., Tesseract or CRNNs) process an image, a robust preprocessing and extraction phase is strictly required.

This paper outlines a modular, offline-capable Deep Learning preprocessing layer written entirely in C++17. By distancing the pipeline from traditional Python-heavy dependencies, our implemented architecture minimizes latency and improves deployment feasibility across varying host specifications. 

\section{Methodology}
The proposed tool accepts an input image and a pre-trained Open Neural Network Exchange (ONNX) format model, generating localized, perspective-corrected, and binarized image patches ready for character-level recognition.

\subsection{Deep Learning Text Detection (DBNet)}
We utilize the Real-time Scene Text Detection with Differentiable Binarization (DBNet) model, loaded via target ONNX inference backends. ONNX Runtime provides hardware-accelerated computation while maintaining cross-platform compatibility. Model operations utilize an Intermediate Representation footprint mapping to ONNX IR version 9, ensuring stable behavior across differing system constraints.

\subsection{Geometry Processing and Polygon Manipulation}
Detected probabilities are refined into concrete bounding geometries. Utilizing Stable Component Analysis alongside the open-source \textit{Clipper} library algorithms, fragmented components of single text lines are merged and dilated logically. This mitigates character drops caused by spatial breaks in DBNet's heatmaps.

\subsection{Perspective Warping and Local Binarization}
Target regions derived from arbitrary textual bounds occasionally conform to non-rectangular polygons. The \texttt{ImageWarp} module calculates homography profiles and performs projective transformations mapping irregular bounds to straight rectangular dimensions. The final localized bounds are converted to stark black and white representations, negating background interference in subsequent OCR logic.

\section{System Architecture and Implementation}
Our framework leverages standard CMake structure natively compatible with GCC, Clang, or MSVC. The project components include:
\begin{itemize}
    \item \textbf{TextDetector:} Encapsulates ONNX Runtime inferencing.
    \item \textbf{GeometryUtils:} Intercepts mathematical derivations and shape clustering.
    \item \textbf{ImageWarp:} Enacts mathematical transformations.
    \item \textbf{PostProcess:} Converts unwarped components to binarized states.
\end{itemize}

\section{Evaluation and Results}
(Author Note: In this section, insert qualitative or quantitative metrics showing how the preprocessing improves OCR accuracy compared to raw image input, along with elapsed CPU/GPU inference times.)

\section{Conclusion}
This work demonstrates a lightweight, portable, and deterministic C++ mechanism to augment state-of-the-art OCR workflows. By combining Differentiable Binarization with intricate C++ geometric manipulation and image unwarping, OCR engine recognition fidelity can be maximized dynamically.

\begin{thebibliography}{00}
\bibitem{b1} M. Liao, Z. Wan, C. Yao, K. Chen, and X. Bai, ``Real-time Scene Text Detection with Differentiable Binarization,'' in Proc. AAAI, 2020.
\bibitem{b2} ONNX Runtime Developers, ``ONNX Runtime: Cross-platform, high performance machine learning inferencing and training accelerator,'' 2021. [Online]. Available: https://onnxruntime.ai/
\bibitem{b3} A. Johnson. \textit{Clipper - an open source freeware polygon clipping library}. [Online]. Available: http://www.angusj.com/delphi/clipper.php
\end{thebibliography}

\end{document}
